\documentclass[../../master-handout.tex]{subfiles}

\begin{document}

\section{Lecture 23 (Wed, Jul 29)}
\sidenote{\textbf{Scribe:} Ehsan Tahvilian (EhsanTahvilian) \\ \textbf{Date:} Wed, Jul 29, 2026}

\subsection{Minimum-norm CLF--CBF control}

This lecture develops an optimization-based controller that combines a stability objective, encoded by a control Lyapunov function (CLF), with a safety or invariance objective, encoded by a control barrier function (CBF). The control input is selected by solving a minimum-norm quadratic program (QP). Because the objective is quadratic and the CLF and CBF constraints are affine in the control input, the optimal controller can be characterized using the Karush--Kuhn--Tucker (KKT) conditions. The lecture first derives the four possible active-set solutions of the unrelaxed CLF--CBF QP, then introduces a slack variable that relaxes the CLF constraint while keeping the CBF constraint hard, and finally applies the construction to a single-integrator robot that must reach a desired point while avoiding a circular obstacle.

\subsection{Why nonlinear control is treated numerically}

Linear systems are important because they are analytically tractable and provide useful local approximations of nonlinear systems. Nevertheless, their qualitative behavior is comparatively restricted. For example, the solution of a time-invariant linear system $\dot{x}=Ax$ is expressed through the matrix exponential, $x(t)=e^{A(t-t_0)}x(t_0)$, and therefore a finite initial state does not escape to infinity in finite time. A linear autonomous system also cannot possess an isolated attracting limit cycle. By contrast, nonlinear systems can exhibit finite-time escape, periodic orbits and limit cycles, multiple equilibria, bifurcations, and trajectories for which no useful closed-form expression exists. Depending on the regularity of the vector field, nonlinear initial-value problems can even lose uniqueness.

These features motivate a numerical viewpoint. Rather than deriving one explicit feedback law that handles every operating condition, an optimization-based controller evaluates the control input online. At every state, the controller solves a small optimization problem whose constraints represent the tasks that the robot must satisfy. In the minimum-norm formulation considered here, the controller can be interpreted as projecting the zero input onto the set of inputs that satisfy the CLF and CBF inequalities.

\subsection{Control-affine dynamics and Lie derivatives}

Consider the nonlinear control-affine system
\begin{equation}
    \dot{x}=f(x)+g(x)u,
    \label{eq:lec23-control-affine}
\end{equation}
where $x\in\mathbb{R}^{n}$ is the state, $u\in\mathbb{R}^{m}$ is the control input, $f:\mathbb{R}^{n}\rightarrow\mathbb{R}^{n}$ is the drift vector field, and $g:\mathbb{R}^{n}\rightarrow\mathbb{R}^{n\times m}$ is the input vector field. The vector fields are assumed to have sufficient regularity for the closed-loop initial-value problem to be well posed.

For a continuously differentiable scalar function $\phi:\mathbb{R}^{n}\rightarrow\mathbb{R}$, the derivative of $\phi$ along the trajectories of \eqref{eq:lec23-control-affine} is
\begin{align}
    \dot{\phi}(x,u)
    &=
    \frac{\partial\phi}{\partial x}(x)\dot{x}
    \nonumber\\
    &=
    \frac{\partial\phi}{\partial x}(x)f(x)
    +
    \frac{\partial\phi}{\partial x}(x)g(x)u
    \nonumber\\
    &=
    L_f\phi(x)+L_g\phi(x)u.
    \label{eq:lec23-lie-derivative}
\end{align}
Here,
\begin{equation}
    L_f\phi(x)
    :=
    \frac{\partial\phi}{\partial x}(x)f(x)
    \in\mathbb{R},
    \qquad
    L_g\phi(x)
    :=
    \frac{\partial\phi}{\partial x}(x)g(x)
    \in\mathbb{R}^{1\times m}
\end{equation}
are the Lie derivatives of $\phi$ along $f$ and $g$, respectively. The important feature of \eqref{eq:lec23-lie-derivative} is that the derivative of $\phi$ is affine in the decision variable $u$.

\sidenote{A class-$\mathcal{K}$ function $\alpha$ is continuous, strictly increasing, and satisfies $\alpha(0)=0$. In a CBF condition, an extended class-$\mathcal{K}$ function is commonly used so that the argument can take values on both sides of zero.}

\subsection{Stability and safety as affine constraints}

\paragraph{Control Lyapunov function constraint.}

Let $V:\mathbb{R}^{n}\rightarrow\mathbb{R}_{\geq 0}$ be continuously differentiable and positive definite with respect to a desired equilibrium or desired set. A CLF condition asks for a control input that makes $V$ decrease. Using a class-$\mathcal{K}$ function $\alpha_1$, the desired inequality is
\begin{equation}
    \dot{V}(x,u)
    =
    L_fV(x)+L_gV(x)u
    \leq
    -\alpha_1\!\left(V(x)\right).
    \label{eq:lec23-clf-condition}
\end{equation}
When \eqref{eq:lec23-clf-condition} is satisfied along the closed-loop trajectory, $V$ is forced toward zero according to a prescribed decay condition. Thus, the CLF constraint represents a stability-like or convergence task.

\paragraph{Control barrier function constraint.}

Let the desired safe set be the zero-superlevel set of a continuously differentiable function $h$:
\begin{equation}
    \mathcal{C}
    :=
    \left\{
        x\in\mathbb{R}^{n}
        \;:\;
        h(x)\geq 0
    \right\}.
    \label{eq:lec23-safe-set}
\end{equation}
A CBF condition is used to render $\mathcal{C}$ forward invariant. Using a class-$\mathcal{K}$ or extended class-$\mathcal{K}$ function $\alpha_2$, the constraint is
\begin{equation}
    \dot{h}(x,u)
    =
    L_fh(x)+L_gh(x)u
    \geq
    -\alpha_2\!\left(h(x)\right).
    \label{eq:lec23-cbf-condition}
\end{equation}
The inequality does not require $h$ to increase everywhere. It only prevents $h$ from decreasing so rapidly that a trajectory starting in $\mathcal{C}$ crosses the boundary $h=0$ and enters the unsafe region. In particular, on the boundary of the safe set, where $h(x)=0$, condition \eqref{eq:lec23-cbf-condition} reduces to $\dot{h}(x,u)\geq 0$, which prevents an outward crossing from the safe set.

\subsection{The unrelaxed minimum-norm CLF--CBF QP}

The CLF and CBF constraints can be imposed simultaneously while selecting the smallest control input:
\begin{equation}
\begin{aligned}
    \underset{u\in\mathbb{R}^{m}}{\operatorname{minimize}}
    \quad&
    \|u\|^{2}
    \\
    \text{subject to}
    \quad&
    L_fV(x)+L_gV(x)u
    \leq
    -\alpha_1\!\left(V(x)\right),
    \\
    &
    L_fh(x)+L_gh(x)u
    \geq
    -\alpha_2\!\left(h(x)\right).
\end{aligned}
\label{eq:lec23-unrelaxed-qp}
\end{equation}
For a fixed state $x$, the quantities that depend only on $x$ are constants in the optimization problem. Define
\begin{align}
    A_V(x)&:=L_gV(x),
    &
    b_V(x)&:=L_fV(x)+\alpha_1\!\left(V(x)\right),
    \label{eq:lec23-Av-bv}
    \\
    A_h(x)&:=L_gh(x),
    &
    b_h(x)&:=L_fh(x)+\alpha_2\!\left(h(x)\right).
    \label{eq:lec23-Ah-bh}
\end{align}
The two task inequalities become
\begin{equation}
    A_Vu+b_V\leq 0,
    \qquad
    -A_hu-b_h\leq 0.
    \label{eq:lec23-affine-constraints}
\end{equation}
The dependence on $x$ is omitted temporarily to simplify the notation.

The feasible input set is the intersection of two half-spaces in $\mathbb{R}^{m}$. The level sets of $\|u\|^2$ are spheres centered at the origin. Therefore, if the feasible set is nonempty, the solution of \eqref{eq:lec23-unrelaxed-qp} is the Euclidean projection of $u=0$ onto the feasible set. The objective is strictly convex, so the optimal input is unique whenever the problem is feasible.

\subsection{KKT characterization of the unrelaxed QP}

Associate the nonnegative Lagrange multipliers $\lambda_1$ and $\lambda_2$ with the CLF and CBF inequalities, respectively. The Lagrangian is
\begin{equation}
    \mathcal{L}(u,\lambda_1,\lambda_2)
    =
    u^\top u
    +
    \lambda_1(A_Vu+b_V)
    +
    \lambda_2(-A_hu-b_h).
    \label{eq:lec23-unrelaxed-lagrangian}
\end{equation}
The stationarity condition is
\begin{equation}
    \frac{\partial\mathcal{L}}{\partial u}
    =
    2u^\top+\lambda_1A_V-\lambda_2A_h
    =
    0,
\end{equation}
or, equivalently,
\begin{equation}
    u^\star
    =
    \frac{1}{2}
    \left(
        -\lambda_1A_V^\top
        +
        \lambda_2A_h^\top
    \right).
    \label{eq:lec23-u-from-lambda}
\end{equation}
The complete KKT conditions are
\begin{align}
    2u^\top+\lambda_1A_V-\lambda_2A_h&=0,
    \label{eq:lec23-kkt-stationarity}
    \\
    A_Vu+b_V&\leq 0,
    \label{eq:lec23-kkt-primal-clf}
    \\
    -A_hu-b_h&\leq 0,
    \label{eq:lec23-kkt-primal-cbf}
    \\
    \lambda_1,\lambda_2&\geq 0,
    \label{eq:lec23-kkt-dual}
    \\
    \lambda_1(A_Vu+b_V)&=0,
    \label{eq:lec23-kkt-comp-clf}
    \\
    \lambda_2(-A_hu-b_h)&=0.
    \label{eq:lec23-kkt-comp-cbf}
\end{align}
Complementary slackness implies that each multiplier is zero unless its associated constraint is active. Consequently, there are four active-set cases.

For compactness, define
\begin{equation}
    a:=\|A_V\|^2,
    \qquad
    b:=\|A_h\|^2,
    \qquad
    c:=A_VA_h^\top=A_hA_V^\top.
    \label{eq:lec23-abc}
\end{equation}

\paragraph{Case 1: neither constraint is active.}

If
\begin{equation}
    \lambda_1=0,
    \qquad
    \lambda_2=0,
\end{equation}
then \eqref{eq:lec23-u-from-lambda} gives
\begin{equation}
    u^\star=0.
    \label{eq:lec23-case1}
\end{equation}
This candidate is feasible precisely when the zero input already satisfies both inequalities:
\begin{equation}
    b_V\leq 0,
    \qquad
    b_h\geq 0.
    \label{eq:lec23-case1-feasibility}
\end{equation}
Geometrically, the origin lies inside the intersection of the CLF and CBF half-spaces, so the projection of the origin is the origin itself.

\paragraph{Case 2: only the CLF constraint is active.}

Suppose
\begin{equation}
    \lambda_1>0,
    \qquad
    \lambda_2=0.
\end{equation}
The CLF equality and stationarity relation give
\begin{align}
    u^\star&=-\frac{1}{2}\lambda_1A_V^\top,
    \\
    A_Vu^\star+b_V&=0.
\end{align}
Substituting the first expression into the second gives
\begin{equation}
    -\frac{1}{2}\lambda_1\|A_V\|^2+b_V=0,
\end{equation}
and therefore
\begin{equation}
    \lambda_1
    =
    \frac{2b_V}{\|A_V\|^2},
    \qquad
    u^\star
    =
    -\frac{b_V}{\|A_V\|^2}A_V^\top.
    \label{eq:lec23-case2}
\end{equation}
A nondegenerate Case~2 requires $b_V>0$, so that $\lambda_1>0$, and the resulting input must satisfy the CBF inequality. In terms of $a$, $b$, and $c$, the inactive CBF condition is
\begin{equation}
    -A_hu^\star-b_h<0
    \quad\Longleftrightarrow\quad
    b_h-\frac{b_Vc}{a}>0.
    \label{eq:lec23-case2-inactive}
\end{equation}
Thus, the origin violates the CLF half-space, and the optimal input is its orthogonal projection onto the CLF boundary while remaining strictly inside the CBF half-space.

\paragraph{Case 3: only the CBF constraint is active.}

Suppose
\begin{equation}
    \lambda_1=0,
    \qquad
    \lambda_2>0.
\end{equation}
Stationarity and the active CBF equality give
\begin{align}
    u^\star&=\frac{1}{2}\lambda_2A_h^\top,
    \\
    -A_hu^\star-b_h&=0.
\end{align}
Therefore,
\begin{equation}
    -\frac{1}{2}\lambda_2\|A_h\|^2-b_h=0,
\end{equation}
which yields
\begin{equation}
    \lambda_2
    =
    -\frac{2b_h}{\|A_h\|^2},
    \qquad
    u^\star
    =
    -\frac{b_h}{\|A_h\|^2}A_h^\top.
    \label{eq:lec23-case3}
\end{equation}
A nondegenerate Case~3 requires $b_h<0$, so that $\lambda_2>0$. The CLF constraint must remain inactive:
\begin{equation}
    A_Vu^\star+b_V<0
    \quad\Longleftrightarrow\quad
    b_V-\frac{b_hc}{b}<0.
    \label{eq:lec23-case3-inactive}
\end{equation}
In this case, the origin violates the CBF half-space, and the optimal input is the orthogonal projection of the origin onto the CBF boundary while still satisfying the CLF inequality.

\paragraph{Case 4: both constraints are active.}

Suppose
\begin{equation}
    \lambda_1>0,
    \qquad
    \lambda_2>0.
\end{equation}
Both affine inequalities hold with equality. Substituting \eqref{eq:lec23-u-from-lambda} into the two active constraints gives
\begin{fullwidth}
\begin{equation}
    \frac{1}{2}
    \begin{bmatrix}
        -\|A_V\|^2 & A_VA_h^\top
        \\
        A_hA_V^\top & -\|A_h\|^2
    \end{bmatrix}
    \begin{bmatrix}
        \lambda_1
        \\
        \lambda_2
    \end{bmatrix}
    =
    \begin{bmatrix}
        -b_V
        \\
        b_h
    \end{bmatrix}.
    \label{eq:lec23-case4-linear-system}
\end{equation}
\end{fullwidth}
Define
\begin{equation}
    \Delta
    :=
    \|A_V\|^2\|A_h\|^2
    -
    \left(A_VA_h^\top\right)^2
    =
    ab-c^2.
    \label{eq:lec23-delta}
\end{equation}
When $\Delta>0$, the two constraint normals are linearly independent, and the multiplier equations have the unique solution
\begin{align}
    \lambda_1
    &=
    \frac{
        2\left(
            b\,b_V-c\,b_h
        \right)
    }{
        \Delta
    },
    \label{eq:lec23-case4-lambda1}
    \\
    \lambda_2
    &=
    \frac{
        2\left(
            c\,b_V-a\,b_h
        \right)
    }{
        \Delta
    }.
    \label{eq:lec23-case4-lambda2}
\end{align}
The optimal input is then
\begin{equation}
    u^\star
    =
    \frac{1}{2}
    \left(
        -\lambda_1A_V^\top
        +
        \lambda_2A_h^\top
    \right).
    \label{eq:lec23-case4-u}
\end{equation}
The candidate belongs to Case~4 only when the expressions in \eqref{eq:lec23-case4-lambda1} and \eqref{eq:lec23-case4-lambda2} are both positive. At the optimum, the input lies at the intersection of the two active hyperplanes.

By the Cauchy--Schwarz inequality, $\Delta\geq 0$. If $\theta$ is the angle between the two constraint normals, then
\begin{equation}
    \Delta
    =
    \|A_V\|^2\|A_h\|^2\sin^2\theta.
\end{equation}
Hence, $\Delta=0$ when the normals are linearly dependent. In that degenerate situation, the two constraints are parallel or antiparallel, and they may be redundant or mutually inconsistent. The explicit Case~4 formula should not be evaluated when $\Delta$ is numerically close to zero; a numerical QP solver handles this situation more reliably.

\subsection{Why the CLF constraint is relaxed}

Although the CLF and CBF inequalities may each be feasible separately, their intersection can be empty. This occurs when the stability request and the safety request demand incompatible instantaneous control actions. In a safety-critical controller, it is natural to preserve the CBF constraint and soften the CLF constraint. The CLF describes performance: temporary loss of the desired convergence rate is usually acceptable. The CBF describes safety: violating it may allow the state to leave the safe set.

Introduce a scalar slack variable $\delta\in\mathbb{R}$ and a penalty $\rho>0$. The relaxed QP is
\begin{equation}
\begin{aligned}
    \underset{u\in\mathbb{R}^{m},\,\delta\in\mathbb{R}}{\operatorname{minimize}}
    \quad&
    \|u\|^2+\rho\delta^2
    \\
    \text{subject to}
    \quad&
    L_fV(x)+L_gV(x)u
    \leq
    -\alpha_1\!\left(V(x)\right)+\delta,
    \\
    &
    L_fh(x)+L_gh(x)u
    \geq
    -\alpha_2\!\left(h(x)\right).
\end{aligned}
\label{eq:lec23-relaxed-qp}
\end{equation}
In the compact notation of \eqref{eq:lec23-Av-bv}--\eqref{eq:lec23-Ah-bh}, this becomes
\begin{equation}
\begin{aligned}
    \underset{u,\delta}{\operatorname{minimize}}
    \quad&
    \|u\|^2+\rho\delta^2
    \\
    \text{subject to}
    \quad&
    A_Vu-\delta+b_V\leq 0,
    \\
    &
    -A_hu-b_h\leq 0.
\end{aligned}
\label{eq:lec23-relaxed-compact}
\end{equation}
No separate constraint $\delta\geq 0$ is required for the derivation below. At the optimum, the KKT conditions imply $\delta=\lambda_1/(2\rho)\geq 0$.

The parameter $\rho$ controls the price of violating the desired CLF decrease. A large $\rho$ discourages relaxation and makes the controller approach the unrelaxed solution whenever the unrelaxed QP is feasible. A smaller $\rho$ permits a larger $\delta$ and allows the controller to sacrifice the CLF convergence rate more readily.

\subsection{KKT characterization of the relaxed QP}

The relaxed Lagrangian is
\begin{equation}
\begin{aligned}
    \mathcal{L}(u,\delta,\lambda_1,\lambda_2)
    ={}&
    u^\top u+\rho\delta^2
    \\
    &+
    \lambda_1(A_Vu-\delta+b_V)
    \\
    &+
    \lambda_2(-A_hu-b_h).
\end{aligned}
\label{eq:lec23-relaxed-lagrangian}
\end{equation}
Stationarity with respect to $u$ and $\delta$ gives
\begin{align}
    2u^\top+\lambda_1A_V-\lambda_2A_h&=0,
    \label{eq:lec23-relaxed-stationarity-u}
    \\
    2\rho\delta-\lambda_1&=0.
    \label{eq:lec23-relaxed-stationarity-delta}
\end{align}
Consequently,
\begin{align}
    u^\star
    &=
    \frac{1}{2}
    \left(
        -\lambda_1A_V^\top
        +
        \lambda_2A_h^\top
    \right),
    \label{eq:lec23-relaxed-u-lambda}
    \\
    \delta^\star
    &=
    \frac{\lambda_1}{2\rho}.
    \label{eq:lec23-relaxed-delta-lambda}
\end{align}
The primal feasibility, dual feasibility, and complementary-slackness conditions are
\begin{align}
    A_Vu-\delta+b_V&\leq 0,
    \\
    -A_hu-b_h&\leq 0,
    \\
    \lambda_1,\lambda_2&\geq 0,
    \\
    \lambda_1(A_Vu-\delta+b_V)&=0,
    \\
    \lambda_2(-A_hu-b_h)&=0.
\end{align}

Define the relaxation-modified CLF coefficient
\begin{equation}
    a_\rho
    :=
    \|A_V\|^2+\frac{1}{\rho}
    =
    a+\frac{1}{\rho}.
    \label{eq:lec23-arho}
\end{equation}
The additional term $1/\rho$ appears because the relaxed CLF boundary depends on both $u$ and $\delta$. Equivalently, after defining the scaled variable $s=\sqrt{\rho}\,\delta$, the objective becomes $\|u\|^2+s^2$, while the CLF half-space has the normal $\left[A_V,\,-1/\sqrt{\rho}\right]$. The squared norm of this augmented normal is exactly $a_\rho$.

\paragraph{Relaxed Case 1: neither constraint is active.}

When $\lambda_1=\lambda_2=0$, stationarity gives
\begin{equation}
    u^\star=0,
    \qquad
    \delta^\star=0.
\end{equation}
This candidate is feasible when $b_V\leq 0$ and $b_h\geq 0$, exactly as in the unrelaxed problem.

\paragraph{Relaxed Case 2: only the CLF constraint is active.}

Let $\lambda_1>0$ and $\lambda_2=0$. Using \eqref{eq:lec23-relaxed-u-lambda} and \eqref{eq:lec23-relaxed-delta-lambda} in the active CLF constraint gives
\begin{equation}
    -\frac{1}{2}\lambda_1\|A_V\|^2
    -
    \frac{\lambda_1}{2\rho}
    +
    b_V
    =
    0.
\end{equation}
Hence,
\begin{align}
    \lambda_1
    &=
    \frac{2b_V}{a_\rho},
    \label{eq:lec23-relaxed-case2-lambda}
    \\
    u^\star
    &=
    -\frac{b_V}{a_\rho}A_V^\top
    =
    -\frac{\rho b_V}{1+\rho\|A_V\|^2}A_V^\top,
    \label{eq:lec23-relaxed-case2-u}
    \\
    \delta^\star
    &=
    \frac{b_V}{\rho a_\rho}
    =
    \frac{b_V}{1+\rho\|A_V\|^2}.
    \label{eq:lec23-relaxed-case2-delta}
\end{align}
A nondegenerate relaxed Case~2 requires $b_V>0$ and
\begin{equation}
    b_h-\frac{b_Vc}{a_\rho}>0,
\end{equation}
so that the CBF constraint remains inactive.

\paragraph{Relaxed Case 3: only the CBF constraint is active.}

Let $\lambda_1=0$ and $\lambda_2>0$. Equation \eqref{eq:lec23-relaxed-delta-lambda} gives $\delta^\star=0$, and the CBF-only solution is unchanged:
\begin{align}
    \lambda_2
    &=
    -\frac{2b_h}{\|A_h\|^2},
    \\
    u^\star
    &=
    -\frac{b_h}{\|A_h\|^2}A_h^\top,
    \\
    \delta^\star
    &=0.
    \label{eq:lec23-relaxed-case3}
\end{align}
The candidate requires $b_h<0$ and
\begin{equation}
    b_V-\frac{b_hc}{b}<0.
\end{equation}
The relaxation has no effect in this case because the CLF constraint is inactive.

\paragraph{Relaxed Case 4: both constraints are active.}

When both multipliers are positive, substituting \eqref{eq:lec23-relaxed-u-lambda} and \eqref{eq:lec23-relaxed-delta-lambda} into both active constraints yields
\begin{fullwidth}
\begin{equation}
    \frac{1}{2}
    \begin{bmatrix}
        -\left(\|A_V\|^2+\rho^{-1}\right)
        &
        A_VA_h^\top
        \\
        A_hA_V^\top
        &
        -\|A_h\|^2
    \end{bmatrix}
    \begin{bmatrix}
        \lambda_1
        \\
        \lambda_2
    \end{bmatrix}
    =
    \begin{bmatrix}
        -b_V
        \\
        b_h
    \end{bmatrix}.
    \label{eq:lec23-relaxed-case4-system}
\end{equation}
\end{fullwidth}
Define
\begin{equation}
    \Delta_\rho
    :=
    a_\rho b-c^2
    =
    \left(
        \|A_V\|^2+\frac{1}{\rho}
    \right)
    \|A_h\|^2
    -
    \left(A_VA_h^\top\right)^2.
    \label{eq:lec23-delta-rho}
\end{equation}
For $\Delta_\rho>0$, the multipliers are
\begin{align}
    \lambda_1
    &=
    \frac{
        2\left(
            b\,b_V-c\,b_h
        \right)
    }{
        \Delta_\rho
    },
    \label{eq:lec23-relaxed-case4-lambda1}
    \\
    \lambda_2
    &=
    \frac{
        2\left(
            c\,b_V-a_\rho b_h
        \right)
    }{
        \Delta_\rho
    }.
    \label{eq:lec23-relaxed-case4-lambda2}
\end{align}
The optimal control and relaxation are
\begin{align}
    u^\star
    &=
    \frac{1}{2}
    \left(
        -\lambda_1A_V^\top
        +
        \lambda_2A_h^\top
    \right),
    \label{eq:lec23-relaxed-case4-u}
    \\
    \delta^\star
    &=
    \frac{\lambda_1}{2\rho}
    =
    \frac{
        b\,b_V-c\,b_h
    }{
        \rho\Delta_\rho
    }.
    \label{eq:lec23-relaxed-case4-delta}
\end{align}
The candidate is valid when $\lambda_1>0$ and $\lambda_2>0$. Because $a_\rho>a$ for every finite $\rho$, the relaxed determinant $\Delta_\rho$ is strictly larger than the unrelaxed determinant $\Delta$ whenever $A_h\neq 0$. The slack variable therefore regularizes the CLF direction in the augmented decision space.

\subsection{Example: reaching a target while avoiding a circular obstacle}

Consider a planar robot with the single-integrator dynamics
\begin{equation}
    \dot{x}=u,
    \qquad
    x,u\in\mathbb{R}^{2}.
    \label{eq:lec23-single-integrator}
\end{equation}
This is a control-affine system with
\begin{equation}
    f(x)=0,
    \qquad
    g(x)=I_2.
\end{equation}
Let $x_d\in\mathbb{R}^{2}$ be the desired point. Let an unsafe circular region have center $x_o\in\mathbb{R}^{2}$ and radius $d>0$. The geometric construction is shown in Figure~\ref{fig:lec23-single-integrator-geometry}. The initial state is denoted by $x_0$, and the instantaneous control input $u$ determines the state velocity.

\begin{marginfigure}
    \centering
    \includegraphics[width=0.9\linewidth]{lectures/week-12/lecture-23-images/clf_cbf_single_integrator_geometry.pdf}
    \caption{Single-integrator navigation task used in the lecture. The state starts at $x_0$, is driven toward the desired point $x_d$, and must remain outside the unsafe disk centered at $x_o$ with radius $d$. The arrow $u$ indicates the instantaneous control direction.}
    \label{fig:lec23-single-integrator-geometry}
\end{marginfigure}

\paragraph{CLF for convergence to the desired point.}

Choose
\begin{equation}
    V(x)
    =
    \frac{1}{2}\|x-x_d\|^2.
    \label{eq:lec23-example-V}
\end{equation}
This function is nonnegative and vanishes only at $x=x_d$. Its gradient is
\begin{equation}
    \frac{\partial V}{\partial x}
    =
    (x-x_d)^\top.
\end{equation}
Because $f(x)=0$ and $g(x)=I_2$,
\begin{align}
    L_fV(x)&=0,
    \\
    L_gV(x)&=(x-x_d)^\top.
\end{align}
Therefore,
\begin{equation}
    A_V(x)
    =
    (x-x_d)^\top.
    \label{eq:lec23-example-Av}
\end{equation}
Choose the linear class-$\mathcal{K}$ function
\begin{equation}
    \alpha_1(s)=\gamma_Vs,
    \qquad
    \gamma_V>0.
\end{equation}
Then
\begin{equation}
    b_V(x)
    =
    \frac{\gamma_V}{2}\|x-x_d\|^2,
    \label{eq:lec23-example-bv}
\end{equation}
and the CLF inequality is
\begin{equation}
    (x-x_d)^\top u
    +
    \frac{\gamma_V}{2}\|x-x_d\|^2
    \leq
    \delta.
    \label{eq:lec23-example-clf}
\end{equation}
When $\delta=0$, this condition enforces exponential-type decay of the squared distance to the target.

\paragraph{CBF for obstacle avoidance.}

The safe set is the exterior of the open disk:
\begin{equation}
    \mathcal{C}
    =
    \left\{
        x\in\mathbb{R}^{2}
        :
        \|x-x_o\|^2-d^2\geq 0
    \right\}.
    \label{eq:lec23-example-safe-set}
\end{equation}
Choose the barrier function
\begin{equation}
    h(x)
    =
    \|x-x_o\|^2-d^2.
    \label{eq:lec23-example-h}
\end{equation}
Its gradient is
\begin{equation}
    \frac{\partial h}{\partial x}
    =
    2(x-x_o)^\top.
\end{equation}
Thus,
\begin{align}
    L_fh(x)&=0,
    \\
    L_gh(x)&=2(x-x_o)^\top,
\end{align}
and
\begin{equation}
    A_h(x)
    =
    2(x-x_o)^\top.
    \label{eq:lec23-example-Ah}
\end{equation}
Choose
\begin{equation}
    \alpha_2(s)=\gamma_hs,
    \qquad
    \gamma_h>0.
\end{equation}
Then
\begin{equation}
    b_h(x)
    =
    \gamma_h
    \left(
        \|x-x_o\|^2-d^2
    \right),
    \label{eq:lec23-example-bh}
\end{equation}
and the CBF inequality becomes
\begin{equation}
    2(x-x_o)^\top u
    +
    \gamma_h
    \left(
        \|x-x_o\|^2-d^2
    \right)
    \geq
    0.
    \label{eq:lec23-example-cbf}
\end{equation}
On the obstacle boundary, $\|x-x_o\|=d$, so \eqref{eq:lec23-example-cbf} reduces to
\begin{equation}
    2(x-x_o)^\top u\geq 0.
\end{equation}
The vector $x-x_o$ is the outward normal to the circle. Therefore, this boundary condition removes any inward component of the velocity and prevents the robot from entering the obstacle.

\paragraph{Relaxed CLF--CBF controller for the example.}

Combining \eqref{eq:lec23-example-clf} and \eqref{eq:lec23-example-cbf} gives
\begin{equation}
\begin{aligned}
    \underset{u\in\mathbb{R}^{2},\,\delta\in\mathbb{R}}{\operatorname{minimize}}
    \quad&
    \|u\|^2+\rho\delta^2
    \\
    \text{subject to}
    \quad&
    (x-x_d)^\top u
    +
    \frac{\gamma_V}{2}\|x-x_d\|^2
    \leq
    \delta,
    \\
    &
    2(x-x_o)^\top u
    +
    \gamma_h
    \left(
        \|x-x_o\|^2-d^2
    \right)
    \geq
    0.
\end{aligned}
\label{eq:lec23-example-relaxed-qp}
\end{equation}
Far from the obstacle, the CBF constraint is normally inactive, and the solution is determined primarily by the CLF. As the robot approaches the obstacle, the CBF constraint may become active. The QP then modifies the target-seeking control by removing the unsafe inward component. If the desired CLF decrease and the hard CBF requirement conflict, the slack variable becomes positive. The robot temporarily gives up part of the desired convergence rate, passes safely around the obstacle, and resumes stronger convergence toward $x_d$ after the safety constraint becomes inactive.

For this example,
\begin{align}
    a(x)
    &=
    \|A_V(x)\|^2
    =
    \|x-x_d\|^2,
    \\
    b(x)
    &=
    \|A_h(x)\|^2
    =
    4\|x-x_o\|^2,
    \\
    c(x)
    &=
    A_V(x)A_h(x)^\top
    =
    2(x-x_d)^\top(x-x_o).
\end{align}
These quantities can be substituted directly into the active-set formulas derived above. In an implementation, however, it is generally preferable to solve \eqref{eq:lec23-example-relaxed-qp} numerically at each control update. The numerical QP remains reliable near changes of active set and near configurations where the constraint normals are almost parallel.

\end{document}
